In a binary system involving a compact star, where the companion star does not
overflow its Roche lobe, the primary source of accretion for the compact star
is the stellar wind from the companion star. This wind-fed accretion is
especially significant in systems that include an early-type star (such as an
O or B star, indicated henceforth as an OB star), alongside a compact object
like a neutron star (NS) in a close orbit.

Consider such a NS-OB star binary system where a neutron star of mass
$M_{\mathrm{NS}} = 2.0\,\mathrm{M}_\odot$ and radius
$R_{\mathrm{NS}} = 11$\,km is orbiting in a circular orbit of radius $a$
around the centre of the OB star with velocity
$v_{\mathrm{orb}} = 1.5 \times 10^5\,\mathrm{m\,s^{-1}}$ (see figure below).
Throughout this problem the mass loss from the OB star is assumed to be
spherically symmetric and its rate is
$\dot{M}_{\mathrm{OB}} = 1.0 \times 10^{-4}\,\mathrm{M}_\odot\,
\mathrm{yr}^{-1}$.

\ioaafig{2025_TH_T07-f1.pdf}

\begin{description}
\item[(T07.1)] The accretion radius, $R_{\mathrm{acc}}$, is defined as the
  maximum distance from the NS at which the stellar wind can be captured by
  the NS. If the stellar wind speed at the orbital distance of the NS is
  $v_{\mathrm{w}} = 3.0 \times 10^6\,\mathrm{m\,s^{-1}}$, find
  $R_{\mathrm{acc}}$ for the above system in km using standard escape velocity
  calculations. \hfill[3\,pt]
\item[(T07.2)] Assuming that all captured material is accreted by the NS,
  estimate the mass accretion rate, $\dot{M}_{\mathrm{acc}}$, from the stellar
  wind onto the NS in units of $\mathrm{M}_\odot\,\mathrm{yr}^{-1}$ if
  $a = 0.5$\,au. Neglect the effects of radiation pressure and finite cooling
  time of the accreting gas. \hfill[3\,pt]
\item[(T07.3)] Now consider the situation where the stellar wind speed at the
  orbital distance $a$ (near the NS) becomes comparable with orbital speed of
  the NS. The mass accretion rate from the stellar wind onto the NS in this
  case would be given by an expression of the form
  $\dot{M}_{\mathrm{acc}} = \dot{M}_{\mathrm{OB}} f(\tan\beta, q)$, where
  $q = M_{\mathrm{NS}}/M_{\mathrm{OB}}$ is the mass ratio of the binary and
  $\beta$ is the angle in the frame of the NS between the wind velocity
  direction and radial direction away from the OB star. Obtain the expression
  for $f(\tan\beta, q)$ assuming $M_{\mathrm{OB}} \gg M_{\mathrm{NS}}$.
  \hfill[6\,pt]
\item[(T07.4)] Consider that the fully ionized material accretes radially and
  is hindered by the strong magnetic field $\vec{B}$ of the NS. The pressure
  exerted by the magnetic field is given by $\frac{B^2}{2\mu_0}$. We shall
  assume that the NS has a dipole magnetic field whose magnitude in the
  equatorial plane varies with the distance $r$ from the NS for
  $r \gg R_{\mathrm{NS}}$ as
  \[ B(r) = B_0\left(\frac{R_{\mathrm{NS}}}{r}\right)^3 \]
  where $B_0$ is the magnetic field at the equator of the NS. Assume that the
  axis of the magnetic dipole aligns with the rotation axis of the NS.
  \begin{description}
  \item[(T07.4a)] Obtain the magnetic pressure, $P_{\mathrm{eq,\,mag}}$, in
    the equatorial plane in terms of $B_0$, $R_{\mathrm{NS}}$, $r$, and other
    suitable constants. \hfill[1\,pt]
  \item[(T07.4b)] The maximum distance where the accretion flow is stopped by
    the magnetic field at the equatorial plane is called the magnetospheric
    radius $R_{\mathrm{m}}$. This flow of matter will exert a pressure due to
    the relative motion between incoming stellar wind and the NS. Obtain an
    approximate expression for the critical magnetic field $B_{0,\mathrm{c}}$
    for which $R_{\mathrm{m}}$ coincides with $R_{\mathrm{acc}}$ and calculate
    its value in Tesla. Magnetic effects are neglected for
    $r > R_{\mathrm{m}}$ and consider $v_{\mathrm{w}} \gg v_{\mathrm{orb}}$.
    \hfill[7\,pt]
  \end{description}
\end{description}
